
\begin{figure}
    \centering
    \begin{pcvstack}[boxed, space=0.4em]
        \procedure[linenumbering=on,codesize=\small]
        {$\K$}{
            \mathbf{A} \randfrom \RR_q^{k \times \ell} \\
            (\mathbf{s}_1,\mathbf{s}_2)
                \randfrom S_\eta^\ell \times S_\eta^k 
                    \hfill \pccomment{\(S_\eta\) - set of polynomials in \(\RR_q\) with mods coefficients in \([-\eta,\eta]\)} \\
            \mathbf{t} \gets \mathbf{A}\mathbf{s}_1 + \mathbf{s}_2 \\
            \pcreturn (\pk=(\mathbf{A},\mathbf{t}),
                \sk=(\mathbf{s}_1,\mathbf{s}_2)) \hfill \pccomment{Computing \((\mathbf{s}_1,\mathbf{s}_2)\) from \((\mathbf{A},\mathbf{t})\) is an instance of MLWE.}
        }

        \procedure[linenumbering=on,codesize=\small]
        {$\sign(\sk,M)$}{
            \mathbf{z} \gets \bot \\
            \pcwhile \mathbf{z}=\bot \pcdo \\
            \pcind \mathbf{y} \randfrom S_{\gamma_1}^\ell \label{step:sample-y}\\
            \pcind \mathbf{w} \gets
                \mathsf{HighBits}(\mathbf{A}\mathbf{y}, 2\gamma_2)  
                    \hfill\pccomment{\(\mathsf{HighBits}\) is defined in Figure~\ref{fig:HiLoBits}.} \\
            \pcind c \in B_\tau \gets H(M \parallel \mathbf{w}_1)  \label{step:hash-M}
                \hfill \pccomment{challenge} \\
            \hfill \pccomment{} \\
            \pcind \mathbf{z} \gets \mathbf{y}+c\mathbf{s}_1 
                \hfill \pccomment{response} \\
            \pcind \pcif
                \|\mathbf{z}\|_\infty \geq \gamma_1-\beta
                \text{ or }
                \|\mathsf{LowBits}( \mathbf{A}\mathbf{y}-c\mathbf{s}_2, 2\gamma_2)\|_\infty \geq \gamma_2-\beta \\
                \hfill\pccomment{\(\mathsf{LowBits}\) is defined in Figure~\ref{fig:HiLoBits}.} \\
                \pcind \pcind \mathbf{z}\gets\bot 
                \hfill\pccomment{rejection sampling} \\ 
            \pcreturn \sigma=(\mathbf{z},c)
        }

        \procedure[linenumbering=on,codesize=\small]
        {$\ver(\pk,M,\sigma=(\mathbf{z},c))$}{
            \mathbf{w}' \gets
                \mathsf{HighBits}( \mathbf{A}\mathbf{z}-c\mathbf{t}, 2\gamma_2)
                    \hfill\pccomment{\(
                        \mathbf{A}\mathbf{z}-c\mathbf{t} = 
                        \mathbf{A}(\mathbf{y}+c\mathbf{s}_1)-c(\mathbf{A}\mathbf{s}_1 + \mathbf{s}_2 ) = 
                        \mathbf{A}\mathbf{y}-c\mathbf{s}_2\)} \label{step:explainLoHi}\\
            \pccomment{The \(\mathsf{LowBits}\) of \(-c\mathbf{s}_2\) needs to be sufficiently small such that
                        \(\mathsf{HighBits}(\mathbf{w}') = \mathsf{HighBits}(\mathbf{w})\).}\\
            \pcreturn
                \left[
                \|\mathbf{z}\|_\infty
                < \gamma_1-\beta
                \right]
                \land
                \left[
                c=H(M\parallel\mathbf{w}_1')
                \right]
        }
    \end{pcvstack}
    \caption{Simplified ML-DSA signature scheme.}
    \label{fig:ml-dsa}
\end{figure}



\begin{figure}
    \centering
    \begin{pcvstack}[boxed, space=0.4em]
        \procedure[linenumbering=on,codesize=\small]
        {$\mathsf{Decompose}(r,\alpha)$}{
            \hfill \pccomment{\(r\in[0,q-1]\) and \(\alpha\) satisfies $q-1=m\alpha$ for some $m$.} \\
            r_H = r \text{mods} \alpha \\
            r_L = (r-r_0)/\alpha \\
            \pcreturn (r_H,r_L)
        }
        \procedure[linenumbering=on,codesize=\small]
        {$\mathsf{HighBits}(r,\alpha)$}{
            (r_H,r_L) \gets \mathsf{Decompose}(r,\alpha) \\
            \pcreturn r_H
        }
        \procedure[linenumbering=on,codesize=\small]
        {$\mathsf{LowBits}(r,\alpha)$}{
            (r_H,r_L) \gets \mathsf{Decompose}(r,\alpha)\\
            \pcreturn r_L
        }
    \end{pcvstack}
    \caption{Decompose, Highbits and Lowbits}
    \label{fig:HiLoBits}
\end{figure}
